\section{预期结果}

在前文已经完成选题重构、方向比较、\texttt{Belief-State Controller} 的问题定义，以及对创新性、难点和时间可行性的分析之后，本节需要进一步明确：如果本课题沿当前路径推进，最终应形成怎样的成果，才能使其在论文、系统和答辩三个层面都构成一个完整的毕业设计闭环。换句话说，本节讨论的不是“理想情况下系统还可以做到什么”，而是“在当前研究边界和时间条件下，本课题合理且必要的最终产出是什么”。

\subsection{预期形成清晰的问题定义与理论叙述}

首先，本课题最基础、也是最重要的预期结果之一，是形成一个足够清晰、能够独立成立的问题定义。与此前将系统表述为“生命科学多智能体助手”不同，本课题需要最终明确地把研究对象定义为：面向高代价科学工作流的运行时治理控制问题。该问题定义应能够清楚说明：工作流为何不是普通问答或单步工具调用，系统为何面临部分可观测环境，难度、风险、恢复空间和人工介入价值为何应被视为控制变量，以及控制器的动作为什么是工作流治理动作而不是普通工具路由动作。

这一理论叙述的价值在于，它将决定整个毕设是否拥有明确的学术边界。如果最终论文仍然只是围绕“系统做了哪些功能”展开，那么即使系统实现较完整，也难以形成研究问题层面的说服力。因此，预期成果的第一层不是代码，而是一个能够在论文和答辩中自洽成立的问题形式化。

\subsection{预期形成一个可解释的 Belief-State Controller 原型}

在此基础上，第二个核心预期结果应当是一个可运行、可解释、可嵌入现有系统的 \texttt{Belief-State Controller} 原型。这里需要强调的是，该原型并不意味着实现一个理论上完备的 POMDP 求解器，也不意味着构造端到端强化学习控制策略。更合理的预期，是形成一个结构化、可解释的控制层，使其能够利用当前系统中的运行时观测，对关键治理动作做出动态判断。

这一控制器原型至少应体现以下特征：它不只是对现有规则的机械封装，而是基于状态、观测和动作之间的明确关系进行决策；它能够将当前系统已有的运行时信息，例如候选分歧、质量门控结果、恢复轨迹和人工介入记录，转化为对当前执行状态的更新依据；它能够在现有工作流中的关键控制点真正影响系统行为，而不是停留在离线分析或结果说明层面。

\subsection{预期在现有蛋白设计工作流上完成验证}

第三个重要的预期结果，是在当前已有的蛋白设计工作流上完成针对 controller 的验证。由于当前系统已经具备较完整的工作流链条、恢复机制和实验基础，因此本课题并不需要重新寻找新的验证场景，而应当坚持在现有蛋白设计流程中验证控制器的有效性。这样做，一方面能够最大限度复用已有系统能力，另一方面也能使课题始终保持在“高代价科学工作流”这一核心语境中，而不至于因扩展过多场景而分散主线。

从验证目标上看，本课题的预期不是证明 controller 在所有指标上都绝对优于一切基线，而是证明它相较于静态策略、固定规则或不带显式状态更新的简单门控方法，能够在若干关键维度上体现更合理的权衡，例如任务推进率、恢复代价、人工介入次数、执行时间以及失败样本的可解释性等。

\subsection{预期形成与基线策略的对照结论}

除了实现控制器本身，另一个不可缺少的预期结果是形成与基线策略的对照分析。若缺乏对照，即使控制器可以运行，也难以说明其价值所在。因此，本课题应至少形成一组清晰的比较结论，用于回答以下问题：如果不引入 \texttt{Belief-State Controller}，而只是依赖原有静态流程、简单阈值门控或固定恢复规则，系统会表现成什么样；而引入 controller 之后，系统在决策方式和运行结果上又发生了哪些变化。

这种对照的重要性不仅在于实验意义，还在于它使课题的研究主线能够被更清晰地看见。因为只有当控制器能够在相同工作流条件下改变系统行为，并带来可观察的结果差异时，它才真正构成“方法”。

\subsection{预期形成可写入论文与答辩的贡献结构}

从毕业设计的最终交付要求来看，系统实现和实验结果本身并不足够，课题还必须能够转化为清晰的论文结构与答辩叙述。因此，本课题的另一个关键预期结果，是形成一组可以明确写进论文和答辩材料的贡献结构。结合前文分析，这些贡献应主要围绕以下几个方面展开：
\begin{itemize}
	\item 将高代价科学工作流重新定义为一个部分可观测的治理控制问题；
	\item 提出一种面向当前场景的结构化隐状态表示；
	\item 提出一种基于运行时异构观测的信念状态更新与动作选择思路；
	\item 在现有蛋白设计工作流中验证该控制方法的有效性。
\end{itemize}

\subsection{预期结果的边界}

在说明预期成果时，还必须主动界定哪些结果并不是本课题的合理目标。首先，本课题不以追求一个理论上最优、可泛化到所有科学工作流场景的控制器为目标；其次，本课题也不以构建一个完整的端到端强化学习框架为目标；再次，本课题不要求同时完成大规模 benchmark、复杂泛化验证以及多场景迁移。

因此，本课题更合理的预期是：在现有系统之上，围绕一个明确的工作流控制问题，形成清晰的问题形式化、一个可解释的控制器原型、一组与基线的对照实验结论，以及一套能够支撑论文和答辩的贡献表达。

\subsection{本节结论}

综合来看，本课题的预期结果可以概括为四个层次：其一，形成一个明确的“高代价科学工作流治理控制”问题定义；其二，在现有系统之上形成一个可解释的 \texttt{Belief-State Controller} 原型；其三，在蛋白设计工作流场景中形成与基线策略的对照验证结果；其四，将上述内容组织成一套清晰的论文与答辩贡献结构。只要这四个层次能够完成，本课题就能够从一个已有工作流平台中的局部增强，真正转化为一个边界清晰、主线明确、兼具算法性与系统落地性的毕业设计成果。
